\section{Motivation and Scope}
Fraud occurrence and platform intervention are different events. A user may
commit fraud before the platform accumulates sufficient evidence to suspend
the account. Learning suspension time as though it were fraud onset therefore
mixes user behavior with the platform's detection mechanism.

SAFE~\cite{zheng2019safe} combines recurrent sequence modeling with positive
hazard outputs and a monotonically decreasing survival score. Its revised
objective rewards event probability accumulated before an observed suspension,
rather than concentrating event probability at suspension time. This supports
early flagging, but does not explicitly identify when fraud began or how long
the platform took to detect it. Moreover, absence of suspension is not evidence
that fraud has not occurred.

The supplied conceptual note~\cite{smola2022note} proposes separate hazards
for latent fraud occurrence and subsequent detection, together with latent
variable inference. We use that proposal as the starting point of SAFEPlus;
the two-process idea itself is prior conceptual material, not a new
contribution claimed here. Attention, hierarchical representations, and future
distillation described in that note are deferred while the likelihood and its
limitations are examined.

\paragraph{Working research question.}
Under which assumptions can delayed detection records support useful
inference about fraud onset, and when does explicitly modeling detection
delay improve early warning relative to SAFE?

\paragraph{Current scope.}
This draft contains a candidate likelihood, a finite-grid implementation
specification, and an evaluation plan. It reports an exploratory SAFE baseline rerun. A potential
submission contribution must be established through a broader literature
review, explicit identification assumptions, and evidence of reliable
inference or improved decisions. Replacing a GRU with attention alone is not
the present research objective.
